\chapter{In-House EO Infrastructure for Education and Research Innovation}

\section{Purpose and scope}

An in-house Earth Observation (EO) environment can make data handling and analytical
assumptions available for inspection during research and teaching. In this thesis, its
value is assessed through capabilities implemented and exercised in the repository, not
through assumed educational outcomes. The objective is to support research and teaching
within Kvarken Space Center and the University of Vaasa with a modest, reproducible
prototype that can run locally and expose its evidence and limits.

\section{Capabilities available for student use}

The current codebase supports a sequence of practical exercises:
\begin{itemize}
\item Offline, fixture-based metadata searches and experiments, including deterministic
  timeout, rate-limit and malformed-input cases.
\item Source validation and transformation of provider-shaped records into typed,
  transport-neutral \texttt{EOScene} values.
\item Inspection of content-addressed source payloads and append-only provenance records.
\item Regional and quality-filtered queries against the persistent SQLite spatial catalog.
\item Read-only query exercises through the educational HTTP API.
\item Local processing and evaluation of the retained Sentinel raster case study when
  optional scientific dependencies and input data are available.
\item A first remote workflow on CSC in which five authentic captured scene-metadata records
  passed through provenance, SQLite catalog persistence, regional queries and a loopback API.
\end{itemize}

The CSC record does not include remote satellite-image classification. Raster classification
and scientific evaluation ran locally. The CDSE adapter is an implemented metadata interface,
but its tests use offline transports and do not provide a live-provider exercise.

\section{Demonstrated capabilities and proposed benefits}

The implemented exercises make processing steps, source records, validation decisions and
query results inspectable. Offline fixtures allow the metadata path to be run without
credentials or external services. Provider-specific interfaces and the transport-neutral
scene model provide a basis for adding other sources, but provider-neutral operation has not
been demonstrated against a range of live providers.

Future teaching work could turn these components into guided labs on metadata quality,
provenance, geospatial filtering, failure handling and reproducible scientific evaluation.
Future research prototypes could add provider adapters or analytical methods while retaining
the offline test path. Migration from local execution to CSC has been demonstrated only for
the five-record metadata workflow; portability of larger raster workloads remains untested.
These are proposed extensions, not evaluated benefits.

\section{Resource requirements and maintenance}

The metadata baseline uses the Python standard library, SQLite and small fixtures. The
optional raster case study requires additional scientific packages, network access to
reacquire its large input windows, local storage, and enough memory for geospatial window
processing. Those input files are not included in Git. The CSC demonstration used one
small CPU instance for five metadata records; it does not establish minimum resources for
other workloads or any cloud cost advantage.

An in-house environment also requires maintenance: dependency and provider-interface updates,
data-license and attribution checks, credential protection, storage and retention management,
backup procedures, documented commands, and review of changes in product processing. The
educational API is read-only and was tested on loopback; public deployment, authentication,
availability and multi-user operations were not evaluated. No claim is made about improved
learning outcomes, student satisfaction, comparative superiority or cost savings. An
educational evaluation would be needed to assess those outcomes.
